\section{Metodyka pomiarów}
\label{sec:metodyka}
\subsection{Część I: kwalifikacja pomieszczenia}
\begin{enumerate}
\item Zmierzyć geometrię pomieszczenia i badanego źródła.
\item Opisać sposób pomiaru czasu pogłosu, punkty, powtórzenia i konfigurację
ustrojów pochłaniających; zestawić czasy w pasmach oktawowych.
\item Dla prostopadłościennej i półsferycznej powierzchni pomiarowej,
przy odległości $d=\qty{1}{\metre}$ wskazanej w konspekcie kwalifikacji,
wyznaczyć pole powierzchni i poprawki środowiskowe metodą pogłosową.
\item Oszacować chłonność i poprawkę metodą przybliżoną; podać źródło
przyjętych współczynników pochłaniania.
\item Ocenić przydatność pomieszczenia; według konspektu do wyznaczenia
$K_{2A}$ przyjmuje się czas pogłosu dla \qty{1}{\kilo\hertz}.
\end{enumerate}

\subsection{Część II: poziom mocy akustycznej}
\begin{enumerate}
\item Dla każdego z trzech położeń źródła określić powierzchnię pomiarową,
jej pole i współrzędne mikrofonów; dołączyć rysunek z numeracją punktów.
\item Zarejestrować poziomy z pracującym źródłem oraz tło przy wyłączonym
źródle, z zachowaniem porównywalnych ustawień i czasów pomiaru.
\item Zestawić wyniki A, liniowe i w pasmach tercjowych.
\item Wyznaczyć średnie powierzchniowe, poprawki tła i środowiska oraz
poziom mocy akustycznej A. Sprawdzić spełnienie kryteriów metody.
\item Porównać wyniki dla środka pomieszczenia, ściany i naroża.
\end{enumerate}
\todo{Opisać faktyczny przebieg i odstępstwa. Konspekt mocy podaje różne
zakresy odległości w opisie i w przebiegu ćwiczenia; wpisać faktycznie
zastosowaną geometrię i uzgodniony wariant.}
